% cjt.tex -- CredalJT (scheme D5): algorithm, worked example, exactness theory.
%
% Standalone: own preamble and notation; compiles on its own with two pdflatex
% passes and no bibtex step (this repository's exactness notes cite external
% results inline in prose rather than through a .bib; see docs/properties.tex).
%
% Exposition follows NeurIPS practice: the algorithm is described
% mathematically, not as an account of the implementation. An implementation
% detail appears only where it changes a stated guarantee (the local-solution
% inner bound, the denominator floor, the uncertified fallback).
%
% The NeurIPS class file is not vendored in this repository, so the document
% class is plain `article', matching the sibling notes docs/properties.tex and
% docs/cjt_exactness.tex.
%
% Macros are declared locally rather than via % Macros for the NeurIPS 2026 LCN paper

% Calligraphic shortcuts
\newcommand{\cL}{\mathcal{L}}
\newcommand{\cS}{\mathcal{S}}
\newcommand{\cF}{\mathcal{F}}
\newcommand{\cG}{\mathcal{G}}
\newcommand{\cN}{\mathcal{N}}
\newcommand{\cA}{\mathcal{A}}
\newcommand{\cK}{\mathcal{K}}
\newcommand{\cT}{\mathcal{T}}
\newcommand{\cO}{\mathcal{O}}

% Bold shortcuts
\newcommand{\bX}{\mathbf{X}}
\newcommand{\bY}{\mathbf{Y}}
\newcommand{\bZ}{\mathbf{Z}}
\newcommand{\bx}{\mathbf{x}}
\newcommand{\by}{\mathbf{y}}
\newcommand{\bz}{\mathbf{z}}
\newcommand{\bl}{\mathbf{l}}
\newcommand{\bu}{\mathbf{u}}
\newcommand{\bq}{\mathbf{q}}
\newcommand{\bp}{\mathbf{p}}
\newcommand{\bm}{\mathbf{m}}
\newcommand{\bo}{\mathbf{o}}

% Operators
\newcommand{\pa}{\mathrm{pa}}
\newcommand{\bd}{\mathrm{bd}}
\newcommand{\scope}{\mathrm{scope}}
\newcommand{\ext}{\mathrm{ext}}
\newcommand{\prune}{\mathrm{prune}}
\newcommand{\VE}{\mathrm{VE}}
: that file
% targets neurips_2026.tex and its \pa/\ext would collide with the definitions
% below.

\documentclass[11pt]{article}
\usepackage[margin=1in]{geometry}
\usepackage{amsmath,amssymb,amsthm}
\usepackage{booktabs}
\usepackage{hyperref}
\usepackage{tikz}
\usepackage{xcolor}
\usepackage{algorithm}
\usepackage{algpseudocode}
\usetikzlibrary{arrows.meta,positioning,fit,backgrounds,shapes.geometric,calc}

% Long \texttt{} file names are unbreakable; allow a little extra interword
% stretch so they do not overflow the text block.
\setlength{\emergencystretch}{3em}

\newtheorem{theorem}{Theorem}
\newtheorem{proposition}{Proposition}
\newtheorem{lemma}{Lemma}
\newtheorem{corollary}{Corollary}
\theoremstyle{definition}
\newtheorem{definition}{Definition}
\newtheorem{example}{Example}
\newtheorem{remark}{Remark}

\newcommand{\lo}{\underline}
\newcommand{\hi}{\overline}
\newcommand{\indep}{\perp\!\!\!\perp}
\newcommand{\DSE}{\mathcal{D}_{\mathrm{SE}}}
\newcommand{\DLCN}{\mathcal{D}_{\mathrm{LCN}}}
\newcommand{\ch}{\mathrm{ch}}
\newcommand{\pa}{\mathrm{pa}}
\newcommand{\atoms}{\mathrm{atoms}}
\newcommand{\scope}{\mathrm{scope}}
\newcommand{\restr}{\!\restriction\!}
\newcommand{\readoff}{\rho}


\title{\textsc{CredalJT}: Exact Marginal Inference in Logical Credal Networks\\
at Treewidth Cost\\
\large Algorithm, a fully worked example, and the exactness theory}
\author{IBM Research}
\date{}

\begin{document}
\maketitle

\begin{abstract}
Marginal inference in a Logical Credal Network (LCN) is the problem of computing
tight lower and upper bounds on the probability of a query atom over the credal
set $\DLCN$ that the program's sentences and independence assumptions induce.
The direct formulation optimizes over the full $2^n$ joint and does not scale;
the credal-network engines scale but optimize over a strictly larger set, the
strong extension $\DSE$, and are therefore loose in ways that are hard to
predict. This note presents \textsc{CredalJT} (scheme~D5), which occupies
neither position: it targets $\DLCN$ \emph{directly}, yet its cost is
exponential in the treewidth $w$ rather than in $n$.

\textsc{CredalJT} builds a single junction tree over the chain-graph moral graph,
assigns each cluster a variable vector over the joint of its atoms, ties adjacent
clusters by separator-consistency equalities, hosts every sentence and every
non-redundant independence equality on a cluster containing its scope, and
optimizes the query read-off with a certified global solver. We give the
algorithm in full (Section~\ref{sec:algo}), write the resulting nonlinear program
formally (Section~\ref{sec:nlp}), work a complete four-atom example end to end
including every one of its $22$ variables and $22$ constraints
(Section~\ref{sec:example}), and consolidate the exactness theory with complete
proofs (Section~\ref{sec:theory}): a repaired running-intersection
reconstruction lemma, the exactness theorem under hypotheses H1--H4, the
unconditional outer-bound guarantee that survives when H2 and H3 fail, and a
corollary giving exactness on the \texttt{ktree-fr} benchmark family at
$O(n2^{k+1})$. Section~\ref{sec:paper} is a self-contained subsection ready for
inclusion in a paper. Every numerical value reported here was computed on the
stated instance and independently cross-checked (Section~\ref{sec:provenance}).
\end{abstract}

\tableofcontents

\section{Introduction}\label{sec:intro}

A Logical Credal Network specifies imprecise probabilistic knowledge by
bounding the probabilities of propositional formulas: a \emph{sentence} is either
$\ell\le P(\varphi)\le u$ (Type-1) or $\ell\le P(\varphi\mid\psi)\le u$
(Type-2), over propositional formulas $\varphi,\psi$ on binary atoms. Together
with the independence assumptions that the Local Markov Condition derives from
the program's structure, a set of sentences determines a closed set $\DLCN$ of
joint distributions, and the inference task is to compute
$\lo P(a)=\min_{p\in\DLCN}P_p(a)$ and $\hi P(a)=\max_{p\in\DLCN}P_p(a)$ for each
atom $a$.

Two families of engines exist. The \emph{direct} formulation writes one variable
per joint state and optimizes over $\DLCN$ itself; it is exact when the solver
certifies global optimality, but its $2^n$ size confines it to small instances,
and even there the certificate sometimes fails to materialize
(Remark~\ref{rem:oracle}). The
\emph{credal-network} family (Credal~VE, Credal~CTE, ApproxLP, Interval~BP,
Credal~IJGP) compiles the program into a chain-graph credal network with interval
local credal sets, and propagates; it scales, but it optimizes over the strong
extension $\DSE\supseteq\DLCN$ and so inherits two distinct looseness gaps
(Section~\ref{sec:prelim-gaps}).

\textsc{CredalJT} belongs to neither family. It does not compile a credal
network: it never reads the per-family interval bounds and never enumerates
extreme points, consuming only the chain graph's \emph{structure} and the LCN's
own sentences and independence assertions. Consequently it optimizes over a
relaxation of $\DLCN$ built from cluster marginals, not over $\DSE$ --- and
under the hypotheses of Theorem~\ref{thm:d5} that relaxation is exact. Its cost
is $\sum_C 2^{|C|}=O(n\,2^{w+1})$, exponential in treewidth rather than in $n$.

\paragraph{Contributions.}
\begin{enumerate}
  \item A complete specification of the algorithm (Section~\ref{sec:algo},
    Algorithm~\ref{alg:cjt}), including the constraint-hosting triage that keeps
    the treewidth small and the query-independence property that lets one
    constraint program serve every atom.
  \item The nonlinear program written formally (Section~\ref{sec:nlp}), with the
    marginalization operator, the cleared Type-2 encoding, and the joint
    (non-decomposed) form of the independence equalities made explicit.
  \item A four-atom example worked end to end (Section~\ref{sec:example}): the
    junction tree, all $22$ decision variables, all $22$ constraints
    transcribed individually, the resulting bounds, and an independent
    cross-check against a certified global reference on every atom.
  \item The exactness theory consolidated with complete proofs
    (Section~\ref{sec:theory}): Lemma~\ref{lem:rip} (running-intersection
    reconstruction, stated correctly in the presence of zero separator cells),
    Theorem~\ref{thm:d5} (exactness under H1--H4, and the unconditional
    outer-bound half), Remark~\ref{rem:denfloor} (a genuine soundness caveat for
    conditional queries), and Corollary~\ref{cor:ktree} (\texttt{ktree-fr}
    exactness at $O(n2^{k+1})$).
  \item A one-page self-contained subsection for direct inclusion in a paper
    (Section~\ref{sec:paper}).
\end{enumerate}

\section{Preliminaries}\label{sec:prelim}

Throughout, $\mathbf V$ is a set of $n$ binary atoms, $\Delta(2^{\mathbf V})$ the
simplex of joint distributions over the $2^n$ truth assignments, and
$p\restr A$ the marginal of $p$ on an atom subset $A\subseteq\mathbf V$. For a
propositional formula $\varphi$ we write $\atoms(\varphi)$ for its atoms and
$\omega\models\varphi$ for satisfaction.

\subsection{The credal set \texorpdfstring{$\DLCN$}{D\_LCN}}\label{sec:prelim-dlcn}

A Type-1 sentence $\ell\le P(\varphi)\le u$ contributes the two linear rows
$\ell\le\sum_{\omega\models\varphi}p(\omega)\le u$. A Type-2 sentence
$\ell\le P(\varphi\mid\psi)\le u$ is read in \emph{cleared} form,
\begin{equation}\label{eq:cleared}
  \ell\cdot P(\psi)\;\le\;P(\varphi\wedge\psi)\;\le\;u\cdot P(\psi),
\end{equation}
which avoids a division and is the form imposed throughout.

\begin{remark}[The cleared reading is the one in force]\label{rem:cleared}
The cleared rows \eqref{eq:cleared} are vacuously satisfied on
$\{p:P(\psi)=0\}$, so they admit joints that the ratio reading
$\ell\le P(\varphi\wedge\psi)/P(\psi)\le u$ excludes. The two readings therefore
define different credal sets. Every statement in this note is made for the
cleared reading.
\end{remark}

The Local Markov Condition attaches to the program a set of conditional
independence assertions $(X\indep Y\mid S)$ with $X$ a singleton. Each is imposed
as the family of \emph{joint} equalities
\begin{equation}\label{eq:lmc}
  P(\mathbf x,\mathbf y,\mathbf s)\,P(\mathbf s)
  =P(\mathbf x,\mathbf s)\,P(\mathbf y,\mathbf s),
\end{equation}
one for each configuration $(\mathbf y,\mathbf s)$ of $Y\cup S$. Writing $g_s$
for the functional of a sentence $s$ and $h_\iota$ for the residual of an
assertion $\iota$,
\begin{equation}\label{eq:dlcn}
  \DLCN=\bigl\{\,p\in\Delta(2^{\mathbf V})\;:\;
    \ell_s\le g_s(p)\le u_s\ \forall s,\qquad h_\iota(p)=0\ \forall \iota\,\bigr\},
\end{equation}
and the inference targets are
$\lo P(a)=\min_{p\in\DLCN}P(a)$, $\hi P(a)=\max_{p\in\DLCN}P(a)$.

The following elementary lemma is the workhorse of Section~\ref{sec:theory}: it
is what licenses imposing a constraint on a single cluster.

\begin{lemma}[Locality]\label{lem:locality}
Every constraint in \eqref{eq:dlcn} depends on $p$ only through its marginal on a
fixed atom set, its \emph{scope}: a sentence $s$ through $p\restr\atoms(s)$, and
an LMC assertion $(X\indep Y\mid S)$ through $p\restr(X\cup Y\cup S)$.
Consequently, if $p$ and $p'$ satisfy $p\restr A=p'\restr A$ for the scope $A$ of
a constraint $c$, then $c$ holds for $p$ iff it holds for $p'$.
\end{lemma}

\begin{proof}
Let $c$ be a constraint with scope $A$. There are three cases.

\emph{Type-1.} $P(\varphi)=\sum_{\omega\models\varphi}p(\omega)$. Since
$\atoms(\varphi)\subseteq A$, the truth value of $\varphi$ at $\omega$ is
determined by $\omega\restr A$; grouping the sum by the value of $\omega\restr A$
gives $P(\varphi)=\sum_{\alpha\models\varphi}(p\restr A)(\alpha)$, where
$\alpha$ ranges over assignments of $A$. This is a function of $p\restr A$ alone.

\emph{Type-2.} By the previous case both $P(\varphi\wedge\psi)$ and $P(\psi)$
are functions of $p\restr A$ with
$A\supseteq\atoms(\varphi)\cup\atoms(\psi)$; the rows \eqref{eq:cleared} are
therefore functions of $p\restr A$.

\emph{LMC.} Each factor in \eqref{eq:lmc} is a probability of a conjunction of
literals over $X\cup Y\cup S$, hence by the Type-1 case a function of
$p\restr(X\cup Y\cup S)$; a product of such functions is again one.

In each case the constraint's truth value is determined by $p\restr A$, which is
precisely the stated consequence.
\end{proof}

\subsection{Chain graph, strong extension, and the two gaps}\label{sec:prelim-gaps}

The LCN pipeline moralizes and simplifies the program's structure graph into a
\emph{chain graph} whose nodes are chain components; an undirected component
collapses into a single \emph{compound} node (rendered with a \texttt{-}-joined
name, e.g.\ \texttt{C-D}). This yields a factorization
\begin{equation}\label{eq:factor}
  P(\mathbf V)=\prod_v P(\ch_v\mid\pa_v),
\end{equation}
one factor per node $v$, where $\ch_v$ is the (possibly multi-atom) chain
component at $v$ and $\pa_v$ its parents. The \emph{family scope} at $v$ is
$\ch_v\cup\pa_v$; root families have $\pa_v=\varnothing$.

The credal-network engines replace each factor by an interval local credal set
$K_{v,\pi}$ per parent configuration $\pi$ and propagate over the \emph{strong
extension}
\begin{equation}\label{eq:se}
  \DSE=\Bigl\{\textstyle\prod_v P_v\;:\;P_v\in K_{v,\pi}
  \text{ chosen independently for each parent configuration }\pi\Bigr\}
  \;\supseteq\;\DLCN .
\end{equation}
The inclusion is generally strict, for two reasons that it is worth naming since
they are exactly what \textsc{CredalJT} avoids:
\begin{itemize}
  \item \textbf{Gap~A (local-set widening).} Each $K_v$ is computed per family in
    isolation, so it may admit conditionals that no global $p\in\DLCN$ realizes:
    it ignores sentences on \emph{other} families and the cross-family LMC
    equalities.
  \item \textbf{Gap~B (strong-extension widening).} Even with exact local sets,
    the free product in \eqref{eq:se} lets each family pick its row
    independently, which can violate a sentence whose scope spans several
    families.
\end{itemize}
\textsc{CredalJT} reads neither the interval bounds nor the extreme points of any
$K_{v,\pi}$; it uses \eqref{eq:factor} only for its \emph{scopes}. Neither gap
therefore arises, and the target is $\DLCN$ itself. This is the
structural reason it is the only engine in the credal-network directory whose
bound character is stated against $\DLCN$ rather than $\DSE$.

\begin{definition}[Bound character]\label{def:char}
Let $\mathcal R$ be a reference set (here $\DLCN$ or $\DSE$) with true bounds
$[\lo q^\star,\hi q^\star]$ for the query at hand. An engine's interval
$[\lo q,\hi q]$ is \emph{exact w.r.t.\ $\mathcal R$} if
$[\lo q,\hi q]=[\lo q^\star,\hi q^\star]$; a valid \emph{outer} bound if
$[\lo q,\hi q]\supseteq[\lo q^\star,\hi q^\star]$; an \emph{inner} bound if
$[\lo q,\hi q]\subseteq[\lo q^\star,\hi q^\star]$; and \emph{unreliable} if
neither containment is guaranteed. An outer bound is sound as a guarantee; an
inner bound is not.
\end{definition}

\section{The \textsc{CredalJT} algorithm}\label{sec:algo}

\subsection{Overview}\label{sec:algo-overview}

\textsc{CredalJT} proceeds in four steps. It builds a cluster tree over the atoms
from the family scopes of the chain graph \eqref{eq:factor}; it assigns every
sentence and every non-redundant independence assertion to a cluster containing
its scope; it writes a single constraint system whose variables are one
distribution per cluster; and it reads the query off a cluster containing it,
minimizing and maximizing over that system. Algorithm~\ref{alg:cjt} states the
procedure; Section~\ref{sec:nlp} writes the constraint system formally.

The only input the algorithm takes from the chain graph is the collection of
family \emph{scopes} $\{\ch_v\cup\pa_v\}_v$, together with the LCN's own
sentences and assertions. It never forms the local credal sets $K_{v,\pi}$ of
\eqref{eq:se}, and never enumerates their extreme points. This is the structural
reason its target is $\DLCN$ rather than $\DSE$: neither Gap~A nor Gap~B can
arise from a construction that does not build the objects in which those gaps
live.

Throughout, $\mathcal C$ denotes the clusters, $A(C)\subseteq\mathbf V$ the atoms
of cluster $C$, $S_{CD}=A(C)\cap A(D)$ the separator on a tree edge, and
\begin{equation}\label{eq:width}
  w+1=\max_{C\in\mathcal C}|A(C)|,
\end{equation}
so $w$ is the treewidth of the triangulated graph and $w+1$ the largest cluster
measured in atoms. Cluster sizes are always reported in atoms below.

\begin{algorithm}[t]
\caption{\textsc{CredalJT}: marginal bounds at treewidth cost}
\label{alg:cjt}
\begin{algorithmic}[1]
\Require sentences $\Sigma$, LMC assertions $\Lambda$, family scopes $\mathcal F$,
  evidence $e$, width budget $b$
\Ensure $[\lo P(a\mid e),\ \hi P(a\mid e)]$ for every atom $a\notin\atoms(e)$
\State $\mathcal S\gets\mathcal F\cup\{\atoms(s):s\in\Sigma\text{ cross-family}\}$
  \Comment{assertions do not augment}
\State $T\gets\textsc{BucketTree}(\mathcal S,\textsc{MinFillOrder}(\mathcal S))$;
  \ $A(C)\gets\bigcup_{v\in C}\atoms(v)$
  \Comment{RIP holds \hfill [H1]}
\State remove subsumed leaves; \ $S_{CD}\gets A(C)\cap A(D)$ for each edge
\ForAll{$s\in\Sigma$}
  \Comment{every sentence must be hosted \hfill [H2]}
  \State host $s$ on a cluster $C$ with $\atoms(s)\subseteq A(C)$, if one exists
\EndFor
\If{$\max_C|A(C)|>b$ \textbf{or} some $s\in\Sigma$ was not hosted}
  \State \Return $\textsc{Direct}(a,e)$ for every atom $a\notin\atoms(e)$
    \Comment{$2^n$, uncertified; Rem.~\ref{rem:oracle}}
\EndIf
\ForAll{$\iota=(X\indep Y\mid S)\in\Lambda$}
  \Comment{three-way triage \hfill [H3]}
  \If{$S$ separates $X$ from $Y$ in $G^\ast$} \State skip $\iota$
    \Comment{entailed by the separators}
  \ElsIf{$X\cup Y\cup S\subseteq A(C)$ for some $C$} \State host $\iota$ on $C$
  \Else \State drop $\iota$ \Comment{bound stays outer, may be loose}
  \EndIf
\EndFor
\State build $\mathcal G$: \ $q_C\in\Delta(2^{|A(C)|})$,
  \ $q_C\restr S_{CD}=q_D\restr S_{CD}$, \ and \eqref{eq:nlp-t1}--\eqref{eq:nlp-lmc}
  on each host \Comment{once}
\ForAll{atoms $a\notin\atoms(e)$}
  \If{$\{a\}\cup\atoms(e)\subseteq A(C_a)$ for some $C_a$}
    \Comment{objective only \hfill [H4]}
    \State $\lo P\gets\min_{q\in\mathcal G}\readoff_a(q)$, \
      $\hi P\gets\max_{q\in\mathcal G}\readoff_a(q)$
      \Comment{\eqref{eq:nlp-obj}, \eqref{eq:nlp-obj-ev}}
  \Else
    \State $[\lo P,\hi P]\gets\textsc{Direct}(a,e)$
  \EndIf
  \State clamp $[\lo P,\hi P]$ to $[0,1]$
\EndFor
\end{algorithmic}
\end{algorithm}

\subsection{Constructing the cluster tree}\label{sec:algo-jt}

The scopes to be covered are the family scopes, augmented by one clique per
\emph{cross-family sentence} --- a sentence whose atoms fit inside no single
family scope. Augmenting forces such a sentence's atoms to co-occur in a cluster,
which is what guarantees it a host in Section~\ref{sec:algo-host}.

Independence assertions are deliberately \emph{not} used to augment. They are
structural, and a cluster tree with the running-intersection property already
enforces them through its separators (Section~\ref{sec:algo-host}); adding them
as cliques would inflate the width for no gain. A Markov chain, for instance,
would collapse from width~$2$ to a single cluster containing all $n$ atoms.

Given the scopes, an elimination order is computed greedily by \emph{min-fill} on
the interaction graph --- the graph in which two nodes are adjacent when they
co-occur in some scope --- eliminating at each step a vertex whose elimination
adds fewest fill-in edges. Buckets are then formed along that order, each scope
placed in the bucket of the first order variable it contains; a bucket's cluster
is the union of its scopes, and its parent is the first later order variable
appearing in the cluster after the bucket variable is removed. This is standard
bucket elimination, so \textbf{the resulting tree has the running-intersection
property by construction}, discharging hypothesis~H1 of
Theorem~\ref{thm:d5} without further argument.

Each cluster's atom set $A(C)$ is the union of the atoms of its nodes; a compound
node contributes all atoms of its chain component at once, so a cluster may be
wider in atoms than in nodes. Finally, any \emph{leaf} whose atom set is contained
in its parent's is removed, repeatedly: such a leaf is redundant, since the parent
already carries that joint, and removing it is safe because a leaf has no children
to reattach. Only subsumed leaves are removed, so the tree is not claimed to be
minimal --- a non-maximal cluster may survive in the interior.

\subsection{Hosting and the assertion triage}\label{sec:algo-host}

Say a constraint is \emph{hosted} on a cluster $C$ when $\scope(c)\subseteq A(C)$;
by Lemma~\ref{lem:locality} its value is then a function of $q_C$ alone. Hosting
each constraint on a single cluster suffices: separator consistency propagates its
effect to the rest of the tree, as the reconstruction argument of
Theorem~\ref{thm:d5} makes precise.

Every sentence is hosted if it can be. A sentence with no host cannot be
represented exactly on this tree, and the algorithm falls back to the direct
formulation (Section~\ref{sec:algo-solve}).

Independence assertions admit a three-way triage, and this is the algorithm's most
consequential step:
\begin{enumerate}
  \item \textbf{Separator-entailed $\Rightarrow$ omitted.} If the conditioning set
    $S$ separates $X$ from $Y$ in the chordal graph $G^\ast$ (Definition~\ref{def:d5}),
    every separator-consistent family of cluster marginals already satisfies
    $(X\indep Y\mid S)$, so imposing it would be redundant. This is what keeps a
    chain at width~$2$.
  \item \textbf{Hosted $\Rightarrow$ imposed} as the bilinear equalities
    \eqref{eq:lmc} on the host cluster.
  \item \textbf{Neither $\Rightarrow$ dropped.}
\end{enumerate}
Case~3 is the algorithm's only source of looseness, and it is exactly the failure
of hypothesis~H3. The bound remains valid --- the projection half of
Theorem~\ref{thm:d5} requires neither H2 nor H3 --- but it can be strictly loose;
Remark~\ref{rem:d5reading} exhibits an instance where it is.

\subsection{Solving}\label{sec:algo-solve}

The constraint system of Step~3 mentions no query, and every atom lies in some
cluster because each family scope is itself a clique. One system therefore serves
all marginals: per atom, only the objective changes, and two optimizations are
solved. This is the sense in which the construction is \emph{query-independent}.

\paragraph{Global solution.}
The system is nonconvex: the assertion equalities \eqref{eq:nlp-lmc} are bilinear,
and under evidence so is the ratio equality \eqref{eq:nlp-obj-ev}. Exactness in
Theorem~\ref{thm:d5} refers to the global optimum, so the bounds reported here
were computed with a spatial branch-and-bound solver at gap tolerance $0$.

\begin{remark}[A local solver yields an inner bound]\label{rem:ipoptinner}
Solved instead by a local method from a finite set of starting points --- keeping
the smallest minimum and largest maximum found --- the procedure returns a value
no more extreme than the true optimum in each direction. The resulting interval is
therefore contained in the exact one: an \emph{inner} bound in the sense of
Definition~\ref{def:char}, and not a guarantee. Global solution is what makes
Theorem~\ref{thm:d5} applicable.
\end{remark}

\paragraph{Fallback.}
Two conditions divert a query to the direct $2^n$ formulation: some cluster
exceeding the width budget $b$, and a sentence or query having no host. The first
is global to the instance; the second may be local to one atom, arising when that
atom and the evidence share no cluster, i.e.\ when H4 fails for it alone.

\begin{remark}[The fallback is exact in formulation, not in certificate]\label{rem:oracle}
The fallback optimizes the true $2^n$ program over $\DLCN$, but by a local method,
so its output carries no global guarantee and can be over-tight. Conversely, on
instances where the $2^n$ formulation is solvable but its global solver fails to
certify some endpoint, the cluster program of Section~\ref{sec:nlp} is often the
more reliable route to a bound, even at $n$ small enough for the joint to be
written down.
\end{remark}

A failed optimization is reported as the corresponding vacuous endpoint ($0$ for a
minimum, $1$ for a maximum): uninformative, never unsound. If every marginal
returns vacuous, the instance or the evidence is likely inconsistent.

\subsection{Complexity}\label{sec:algo-complexity}

\emph{Variables.} One per cluster state, so by \eqref{eq:width}
\begin{equation}\label{eq:cost}
  \sum_{C\in\mathcal C}2^{|A(C)|}\;=\;O\!\bigl(n\,2^{w+1}\bigr),
\end{equation}
since $|\mathcal C|=O(n)$, plus one auxiliary scalar per atom under evidence.
Against $2^n$ for the direct formulation, this is the algorithm's reason for
being: the cost is exponential in the treewidth, not in the number of atoms.

\emph{Constraints.} There are $|\mathcal C|$ normalization equalities;
$\sum_{(C,D),\,S_{CD}\ne\varnothing}2^{|S_{CD}|}=O(n2^{w})$ separator equalities;
two rows per hosted sentence; and, per hosted assertion $(X\indep Y\mid S)$,
\begin{equation}\label{eq:lmcrows}
  2^{|Y|}\cdot\max\bigl(1,2^{|S|}\bigr)\quad\text{equalities.}
\end{equation}

\begin{remark}[Which rows are nonconvex, and why the width bound survives]\label{rem:counting}
(a) Only the assertion equalities \eqref{eq:nlp-lmc} and, under evidence, the
ratio equality \eqref{eq:nlp-obj-ev} are nonconvex. The cleared Type-2 rows
\eqref{eq:nlp-t2} are \emph{linear} in $q$, since $\ell$ and $u$ are constants.
Consequently an instance with no evidence and no hosted assertion yields a
\emph{linear} program --- which is exactly the situation in
Section~\ref{sec:example}, despite that instance being loopy.
(b) The count \eqref{eq:lmcrows} appears unbounded in $|Y|$, but a \emph{hosted}
assertion satisfies $X\cup Y\cup S\subseteq A(C)$, so it contributes at most
$2^{w+1}$ rows and \eqref{eq:cost} is unaffected. An assertion with a wide $Y$ is
precisely one that the triage of Section~\ref{sec:algo-host} either discharges
structurally or drops.
\end{remark}

Per atom the cost is two optimizations over the shared system, hence $2|\mathbf V|$
in total.

\section{The nonlinear program}\label{sec:nlp}

Fix a cluster $C$ and index the $2^{|A(C)|}$ assignments of $A(C)$ by
$i$, writing $\omega^C_i$ for the $i$-th. For a formula $\varphi$ with
$\atoms(\varphi)\subseteq A(C)$ let $A^{(C)}_\varphi\in\{0,1\}^{2^{|A(C)|}}$ be
its indicator vector, $A^{(C)}_\varphi[i]=\mathbf 1[\omega^C_i\models\varphi]$,
and abbreviate $\langle v,q_C\rangle=\sum_i v_iq_{C,i}$; thus
$\langle A^{(C)}_\varphi,q_C\rangle=P(\varphi)$ whenever $q_C$ is the marginal of
a joint on $A(C)$.

\paragraph{Variables and normalization.}
One distribution per cluster,
\begin{equation}\label{eq:nlp-var}
  q_C\in[0,1]^{2^{|A(C)|}},\qquad
  \sum_{i=0}^{2^{|A(C)|}-1}q_{C,i}=1
  \qquad\forall C\in\mathcal C .
\end{equation}
There are no separator variables, no per-family conditional variables and no
global joint variable; the tree structure enters only through the equalities
below.

\paragraph{Marginalization.}
Let $\pi_{C\to S}$ map a state of $C$ to its restriction to $S\subseteq A(C)$,
and define the $0/1$ aggregation matrix
\begin{equation}\label{eq:nlp-M}
  M^{C\to S}_{r,j}=\mathbf 1\bigl[\pi_{C\to S}(j)=r\bigr],
  \qquad\text{so}\qquad
  \bigl(M^{C\to S}q_C\bigr)[r]=(q_C\restr S)(r).
\end{equation}

\paragraph{Separator consistency.}
For every edge with a non-empty separator, the two endpoints must agree on it:
\begin{equation}\label{eq:nlp-sep}
  \sum_{j:\,\pi_{C\to S}(j)=r}q_{C,j}
  \;=\;\sum_{k:\,\pi_{D\to S}(k)=r}q_{D,k}
  \qquad\forall(C,D)\in E,\ \forall r\in\{0,\dots,2^{|S_{CD}|}-1\},
\end{equation}
equivalently $q_C\restr S_{CD}=q_D\restr S_{CD}$: linear equalities,
$2^{|S_{CD}|}$ per edge. No messages are computed; the consistency conditions are
imposed simultaneously with all others, which is what distinguishes this
formulation from message passing.

\paragraph{Sentences.}
A Type-1 sentence hosted on $C_h\supseteq\atoms(\varphi)$ contributes
\begin{equation}\label{eq:nlp-t1}
  \ell\;\le\;\bigl\langle A^{(C_h)}_{\varphi},\,q_{C_h}\bigr\rangle\;\le\;u,
\end{equation}
and a Type-2 sentence contributes the cleared pair
\begin{equation}\label{eq:nlp-t2}
  \ell\,\bigl\langle A^{(C_h)}_{\psi},q_{C_h}\bigr\rangle
  \;\le\;\bigl\langle A^{(C_h)}_{\varphi\wedge\psi},q_{C_h}\bigr\rangle
  \;\le\;u\,\bigl\langle A^{(C_h)}_{\psi},q_{C_h}\bigr\rangle .
\end{equation}
No division appears, and consistently with Remark~\ref{rem:cleared} the rows
\eqref{eq:nlp-t2} are vacuously satisfied when $P(\psi)=0$: a Type-2 sentence
constrains nothing on a null conditioning event.

\paragraph{Independence assertions.}
For a hosted $(X\indep Y\mid S)$ with $X=\{x\}$, one equality per configuration
$(\mathbf y,\mathbf s)$. With $S\ne\varnothing$, writing the four indicators for
$\{x{=}1,Y{=}\mathbf y,S{=}\mathbf s\}$, $\{S{=}\mathbf s\}$,
$\{x{=}1,S{=}\mathbf s\}$, $\{Y{=}\mathbf y,S{=}\mathbf s\}$ as
$A_a,A_b,A_c,A_d$:
\begin{equation}\label{eq:nlp-lmc}
  \langle A_a,q_{C_h}\rangle\,\langle A_b,q_{C_h}\rangle
  -\langle A_c,q_{C_h}\rangle\,\langle A_d,q_{C_h}\rangle=0 ,
\end{equation}
which is \eqref{eq:lmc} verbatim; and with $S=\varnothing$ the marginal form
$\langle A_a,q\rangle-\langle A_b,q\rangle\langle A_c,q\rangle=0$ for
$P(x,\mathbf y)=P(x)P(\mathbf y)$. These rows are bilinear, hence nonconvex, and
are the reason a global solver is required.

Two features of this encoding matter. First, only the $x{=}1$ slice is needed:
with $(\mathbf y,\mathbf s)$ fixed a binary $X$ has one degree of freedom, so the
$x{=}0$ equality follows. Second, the configurations $\mathbf y$ must range over
\emph{whole} assignments of $Y$. Replacing $(X\indep Y\mid S)$ by the pairwise
family $\{(X\indep t\mid S)\}_{t\in Y}$ is strictly weaker for $|Y|\ge2$ ---
pairwise independence does not imply joint independence --- and under-constrains
the program, yielding spuriously loose bounds.

\paragraph{Objective.}\label{sec:nlp-obj}
Let $C_a$ be a cluster containing the query atom $a$ and all evidence atoms
(hypothesis H4). Without evidence the read-off is linear,
\begin{equation}\label{eq:nlp-obj}
  \readoff_a(q)=\bigl\langle A^{(C_a)}_{a=1},\,q_{C_a}\bigr\rangle=P(a{=}1).
\end{equation}
With evidence $e$ the ratio is reformulated with one auxiliary $t\in[0,1]$,
\begin{equation}\label{eq:nlp-obj-ev}
  \readoff_a(q)=t,\qquad
  t\cdot\bigl\langle A^{(C_a)}_{e},q_{C_a}\bigr\rangle
  =\bigl\langle A^{(C_a)}_{a=1}\!\cdot\!A^{(C_a)}_{e},\,q_{C_a}\bigr\rangle ,
\end{equation}
so $t=P(a{=}1\mid e)$ whenever $P(e)>0$. Since $t$ is unconstrained where
$P(e)=0$, a global solver must be prevented from exploiting that degeneracy, and a
denominator floor
\begin{equation}\label{eq:nlp-floor}
  \bigl\langle A^{(C_a)}_{e},q_{C_a}\bigr\rangle\;\ge\;\varepsilon
\end{equation}
is imposed for a small $\varepsilon>0$. Remark~\ref{rem:denfloor} shows this
floor is a soundness hypothesis, not a numerical convenience.

\paragraph{The program.}
Collecting \eqref{eq:nlp-var}--\eqref{eq:nlp-obj-ev},
\begin{equation}\label{eq:nlp-all}
\begin{aligned}
\lo P(a),\ \hi P(a)\;=\;\min/\max\quad &\readoff_a(q)\\
\text{s.t.}\quad
&\textstyle\sum_i q_{C,i}=1,\quad q_{C,i}\in[0,1] && \forall C\in\mathcal C\\
&q_C\restr S_{CD}=q_D\restr S_{CD} && \forall (C,D)\in E,\ S_{CD}\ne\varnothing\\
&\text{\eqref{eq:nlp-t1}, \eqref{eq:nlp-t2}} && \forall\text{ hosted sentence}\\
&\text{\eqref{eq:nlp-lmc}} && \forall\text{ hosted assertion.}
\end{aligned}
\end{equation}

\section{A worked example}\label{sec:example}

We work \texttt{examples/ktree\_fr\_k2\_n4\_example.lcn} completely. It is a
$k$-tree with $k=2$ on four atoms --- \emph{loopy} after moralization, so the
example is not a degenerate singly-connected case --- yet small enough that every
variable and every constraint can be written down.

\subsection{The instance}

\begin{align*}
s_0:&\quad 0.084382 \le P(x_1) \le 0.684382\\
s_1:&\quad 0.01 \le P(\lnot x_2 \mid \lnot x_1) \le 0.572656\\
s_2:&\quad 0.092785 \le P(x_3 \mid x_1\wedge x_2) \le 0.692785\\
s_3:&\quad 0.068242 \le P(x_0 \mid x_1\wedge x_3) \le 0.668242
\end{align*}
The chain-graph families, in construction order, are
\[
  x_1\mid\varnothing,\qquad x_2\mid x_1,\qquad x_3\mid x_1,x_2,\qquad
  x_0\mid x_1,x_3,
\]
and the Local Markov Condition yields exactly one assertion,
\[
  (x_0\indep x_2\mid x_1,x_3).
\]

\subsection{The junction tree}

Min-fill elimination produces four clusters of at most three atoms, arranged in a
path (Figure~\ref{fig:jt}). In the convention \eqref{eq:width} this is
$w+1=3$, i.e.\ treewidth $w=2=k$, as expected for a $k$-tree:
\[
  \underbrace{\{x_3\}}_{C_4}\ -\
  \underbrace{\{x_2,x_3\}}_{C_3}\ -\
  \underbrace{\{x_1,x_2,x_3\}}_{C_2}\ -\
  \underbrace{\{x_1,x_3,x_0\}}_{C_1},
\]
with separators $S_{12}=\{x_1,x_3\}$, $S_{23}=\{x_2,x_3\}$, $S_{34}=\{x_3\}$.
Hosting assigns $s_0,s_3$ to $C_1$ and $s_1,s_2$ to $C_2$.

\begin{figure}[ht]\centering
\begin{tikzpicture}[
  atom/.style={circle,draw,minimum size=8mm,inner sep=0pt,font=\small},
  clus/.style={rectangle,draw,rounded corners=2pt,minimum height=8mm,
               minimum width=26mm,inner sep=4pt,font=\small,fill=blue!5},
  sep/.style={rectangle,draw,dashed,minimum height=6mm,minimum width=15mm,
              inner sep=2pt,font=\scriptsize,fill=gray!10},
  panel/.style={font=\small}]

% ---- panel (a): moral graph ----
\begin{scope}[shift={(0,0)}]
  \node[atom] (x0) at (0,3.3)    {$x_0$};
  \node[atom] (x1) at (-1.0,1.9) {$x_1$};
  \node[atom] (x3) at (1.0,1.9)  {$x_3$};
  \node[atom] (x2) at (0,0.5)    {$x_2$};
  \draw (x0)--(x1); \draw (x0)--(x3); \draw (x1)--(x3);
  \draw (x1)--(x2); \draw (x2)--(x3);
  \node[panel] at (0,-0.55) {(a) moral graph (loopy)};
\end{scope}

% ---- panel (b): junction tree, separators as on-edge nodes ----
\begin{scope}[shift={(7.4,0)}]
  \node[clus] (c4) at (0,4.05)  {$C_4{:}\ \{x_3\}$};
  \node[sep]  (s34) at (0,3.05) {$\{x_3\}$};
  \node[clus] (c3) at (0,2.05)  {$C_3{:}\ \{x_2,x_3\}$};
  \node[sep]  (s23) at (0,1.05) {$\{x_2,x_3\}$};
  \node[clus] (c2) at (0,0.05)  {$C_2{:}\ \{x_1,x_2,x_3\}$};
  \node[sep]  (s12) at (0,-0.95){$\{x_1,x_3\}$};
  \node[clus] (c1) at (0,-1.95) {$C_1{:}\ \{x_1,x_3,x_0\}$};
  \draw (c4)--(s34); \draw (s34)--(c3);
  \draw (c3)--(s23); \draw (s23)--(c2);
  \draw (c2)--(s12); \draw (s12)--(c1);
  \node[panel,align=center] at (0,-2.95)
    {(b) cluster tree; separators\\ dashed, on the edges};
\end{scope}
\end{tikzpicture}
\caption{The example instance. Moralization makes $\{x_1,x_2,x_3\}$ and
$\{x_0,x_1,x_3\}$ triangles, so the graph is loopy; the cluster tree nonetheless
has width $2$. The separator $S_{12}=\{x_1,x_3\}$ equals the conditioning set of
the sole assertion $(x_0\indep x_2\mid x_1,x_3)$, which is therefore
separator-entailed and never imposed.}
\label{fig:jt}
\end{figure}

\subsection{Separator-entailment discharges the assertion}

The chordal graph $G^\ast$ of Definition~\ref{def:d5} --- two atoms adjacent iff
they co-occur in a cluster --- has edge set
\[
  x_0x_1,\quad x_0x_3,\quad x_1x_2,\quad x_1x_3,\quad x_2x_3,
\]
and in particular $x_0x_2\notin G^\ast$. Deleting $S=\{x_1,x_3\}$ therefore
disconnects $x_0$ from $x_2$: the only paths between them run through $x_1$ or
$x_3$. So $S$ separates $x_0$ from $x_2$ in $G^\ast$, and the sole assertion
$(x_0\indep x_2\mid x_1,x_3)$ is separator-entailed --- case~1 of the triage in
Section~\ref{sec:algo-host}. It is \emph{not} imposed, and the constraint count
below shows zero assertion rows. By Remark~\ref{rem:counting}(a) the program is
consequently a \emph{linear} program, even though the instance is loopy.

This is the mechanism of Theorem~\ref{thm:d5} in its clearest form: the
independence is enforced by the tree's \emph{topology} --- through
Lemma~\ref{lem:rip}'s reconstruction, which produces a joint factorizing over
the chordal graph and hence satisfying its global Markov property --- rather than
by any algebraic row.

\subsection{The program in full}

\paragraph{Variables ($22$).}
$q_{C_1}\in\Delta(8)$ over $(x_1,x_3,x_0)$; $q_{C_2}\in\Delta(8)$ over
$(x_1,x_2,x_3)$; $q_{C_3}\in\Delta(4)$ over $(x_2,x_3)$;
$q_{C_4}\in\Delta(2)$ over $(x_3)$. Total $8+8+4+2=22$.

\paragraph{Simplex ($4$).}
$\sum_i q_{C_k,i}=1$ for $k=1,\dots,4$.

\paragraph{Separator consistency ($4+4+2=10$).}
Writing $q_{C}[\,\cdot\,]$ for a cell of cluster $C$, edge $C_1\!-\!C_2$ over
$\{x_1,x_3\}$ gives four equalities,
\[
\begin{aligned}
q_{C_1}[x_1{=}0,x_3{=}0,x_0{=}0]+q_{C_1}[x_1{=}0,x_3{=}0,x_0{=}1]
 &= q_{C_2}[x_1{=}0,x_2{=}0,x_3{=}0]+q_{C_2}[x_1{=}0,x_2{=}1,x_3{=}0],\\
q_{C_1}[x_1{=}0,x_3{=}1,x_0{=}0]+q_{C_1}[x_1{=}0,x_3{=}1,x_0{=}1]
 &= q_{C_2}[x_1{=}0,x_2{=}0,x_3{=}1]+q_{C_2}[x_1{=}0,x_2{=}1,x_3{=}1],\\
q_{C_1}[x_1{=}1,x_3{=}0,x_0{=}0]+q_{C_1}[x_1{=}1,x_3{=}0,x_0{=}1]
 &= q_{C_2}[x_1{=}1,x_2{=}0,x_3{=}0]+q_{C_2}[x_1{=}1,x_2{=}1,x_3{=}0],\\
q_{C_1}[x_1{=}1,x_3{=}1,x_0{=}0]+q_{C_1}[x_1{=}1,x_3{=}1,x_0{=}1]
 &= q_{C_2}[x_1{=}1,x_2{=}0,x_3{=}1]+q_{C_2}[x_1{=}1,x_2{=}1,x_3{=}1];
\end{aligned}
\]
edge $C_2\!-\!C_3$ over $\{x_2,x_3\}$ gives four more,
\[
  q_{C_2}[x_1{=}0,x_2{=}\beta,x_3{=}\gamma]+q_{C_2}[x_1{=}1,x_2{=}\beta,x_3{=}\gamma]
  = q_{C_3}[x_2{=}\beta,x_3{=}\gamma],
  \qquad \beta,\gamma\in\{0,1\};
\]
and edge $C_3\!-\!C_4$ over $\{x_3\}$ gives two,
\[
  q_{C_3}[x_2{=}0,x_3{=}\gamma]+q_{C_3}[x_2{=}1,x_3{=}\gamma]
  = q_{C_4}[x_3{=}\gamma],\qquad \gamma\in\{0,1\}.
\]

\paragraph{Sentence rows ($8$).}
On $C_1$, from $s_0$ (Type-1) and $s_3$ (Type-2, cleared per
\eqref{eq:nlp-t2}):
\[
\begin{aligned}
0.084382&\le\bigl\langle A^{(C_1)}_{x_1},q_{C_1}\bigr\rangle\le 0.684382,\\
0.068242\,\bigl\langle A^{(C_1)}_{x_1\wedge x_3},q_{C_1}\bigr\rangle
&\le\bigl\langle A^{(C_1)}_{x_0\wedge x_1\wedge x_3},q_{C_1}\bigr\rangle
\le 0.668242\,\bigl\langle A^{(C_1)}_{x_1\wedge x_3},q_{C_1}\bigr\rangle .
\end{aligned}
\]
On $C_2$, from $s_1$ and $s_2$ (both Type-2):
\[
\begin{aligned}
0.01\,\bigl\langle A^{(C_2)}_{\lnot x_1},q_{C_2}\bigr\rangle
&\le\bigl\langle A^{(C_2)}_{\lnot x_2\wedge\lnot x_1},q_{C_2}\bigr\rangle
\le 0.572656\,\bigl\langle A^{(C_2)}_{\lnot x_1},q_{C_2}\bigr\rangle,\\
0.092785\,\bigl\langle A^{(C_2)}_{x_1\wedge x_2},q_{C_2}\bigr\rangle
&\le\bigl\langle A^{(C_2)}_{x_3\wedge x_1\wedge x_2},q_{C_2}\bigr\rangle
\le 0.692785\,\bigl\langle A^{(C_2)}_{x_1\wedge x_2},q_{C_2}\bigr\rangle .
\end{aligned}
\]

\paragraph{Independence rows ($0$).}
None, by separator-entailment as argued above.

\paragraph{Totals.}
$22$ variables and $22$ constraints ($4$ simplex, $10$ separator, $8$ sentence,
$0$ LMC), against a joint formulation's $2^4=16$ variables. The example is
deliberately small enough that the treewidth formulation is not yet a saving; the
saving is asymptotic, $O(n2^{w+1})$ against $2^n$, and Table~\ref{tab:contrast}
shows the same machinery on two five-atom instances.

\subsection{Results}

Solving \eqref{eq:nlp-all} for each atom with SCIP (global, gap $0$) and
comparing against the independent $2^n$ global engine:

\begin{center}\small
\begin{tabular}{@{}lcccc@{}}
\toprule
engine & $P(x_0)$ & $P(x_1)$ & $P(x_2)$ & $P(x_3)$ \\
\midrule
direct $2^n$ formulation, global & $[0,1]$ & $[0.084382,0.684382]$
  & $[0.134877,0.996844]$ & $[0,1]$ \\
\textsc{CredalJT} (D5, SCIP)              & $[0,1]$ & $[0.084382,0.684382]$
  & $[0.134877,0.996844]$ & $[0,1]$ \\
\bottomrule
\end{tabular}
\end{center}

\noindent The two engines agree on every atom to every digit reported, with all
eight of the reference optimizations certified at optimality gap $0$. \textsc{CredalJT} reports \texttt{d5\_exact} true, maximum
cluster $3$ atoms, and a total time of about $1.8$\,s.

The vacuous $[0,1]$ on $x_0$ and $x_3$ is genuine rather than a solver artifact,
and the independent oracle confirms it: no sentence constrains $P(x_3)$ except
through conditioning events that can themselves be driven to probability zero,
at which point \eqref{eq:nlp-t2} becomes vacuous (Remark~\ref{rem:cleared}) and
frees the atom. Setting $P(x_1\wedge x_3)=0$ likewise releases $s_3$ and hence
$x_0$. Both endpoints are thus attained by feasible witnesses.

The separator marginals attained at an optimum make the mechanism concrete.
Table~\ref{tab:msgs} lists them at the maximizer of $P(x_1)$. Both endpoints of
every edge agree on these values by \eqref{eq:nlp-sep}; this is the calibrated
family of cluster marginals that Lemma~\ref{lem:rip} reassembles into a joint of
$\DLCN$ attaining the bound.

\begin{table}[ht]\centering\small
\caption{Separator marginals $q_C\restr S$ at the maximizer of $P(x_1)$, listed
over the assignments of $S$. Both clusters incident to an edge realize the same
vector, by \eqref{eq:nlp-sep}.}\label{tab:msgs}
\begin{tabular}{@{}llcccc@{}}
\toprule
edge & separator $S$ & \multicolumn{4}{c}{$q\restr S$ by assignment of $S$}\\
\midrule
$C_1\!-\!C_2$ & $\{x_1,x_3\}$ & $00{:}\,0.3125$ & $01{:}\,0.0032$
  & $10{:}\,0.6844$ & $11{:}\,0.0000$\\
$C_2\!-\!C_3$ & $\{x_2,x_3\}$ & $00{:}\,0.6844$ & $01{:}\,0.0032$
  & $10{:}\,0.3125$ & $11{:}\,0.0000$\\
$C_3\!-\!C_4$ & $\{x_3\}$     & $0{:}\,0.9968$  & $1{:}\,0.0032$
  & & \\
\bottomrule
\end{tabular}
\end{table}

\subsection{Two contrasting structures}

The running example imposes no assertion row. Table~\ref{tab:contrast} reports two
five-atom instances that do, so that the bilinear rows are exercised rather than
discharged; one of them also exhibits a compound node.

\begin{table}[ht]\centering\small
\caption{Program size on two further instances, measured. \texttt{asia} has a
hosted \emph{conditional} assertion $(C\indep B\mid S)$ contributing
$2^{|Y|}2^{|S|}=4$ bilinear rows; \texttt{alarm} has a hosted \emph{marginal}
assertion $(B\indep E)$ with $S=\varnothing$ contributing $2^{|Y|}=2$ rows, and
its undirected component becomes the compound node \texttt{C-D} whose cluster
carries atoms $\{A,C,D\}$.}\label{tab:contrast}
\begin{tabular}{@{}lcccrrrrrr@{}}
\toprule
instance & clus. & max $|A(C)|$ & edges & vars & constr. & simplex & sep. &
  sent. & LMC \\
\midrule
\texttt{examples/asia.lcn}  & 5 & 3 & 4 & 30 & 33 & 5 & 14 & 10 & 4 \\
\texttt{examples/alarm.lcn} & 4 & 3 & 3 & 22 & 24 & 4 &  8 & 10 & 2 \\
\bottomrule
\end{tabular}
\end{table}

For \texttt{asia} the tree is the path
$\{X\}-\{D,X\}-\{C,D,X\}-\{B,C,D\}-\{B,S,C\}$ with separators
$\{X\},\{D,X\},\{C,D\},\{B,C\}$, so the separator count is
$2+4+4+4=14$, as \eqref{eq:nlp-sep} predicts. For \texttt{alarm} the tree
branches: $\{E\}-\{E,A\}$ with children $\{B,E,A\}$ and $\{A,C,D\}$, giving
$2+4+2=8$.

\section{Exactness}\label{sec:theory}

This section states and proves the guarantees. The development is closed: it uses
only Lemma~\ref{lem:locality} from Section~\ref{sec:prelim} and
Lemma~\ref{lem:rip} below, and in particular needs no notion of realizability, no
property of the strong extension, and no result about any other engine.

\begin{definition}[Cluster-feasible set, hosting, separator-entailment, read-off]\label{def:d5}
Let $T$ be a clique tree over clusters $\mathcal C$, built by the standard
moralize~$\to$~triangulate~$\to$~clique-tree pipeline, and let $G^\ast$ be the
chordal graph whose maximal cliques are the clusters (two atoms adjacent iff they
co-occur in a cluster).
\begin{itemize}
  \item The \emph{cluster-feasible set} is
  \begin{align*}
    \mathcal G=\bigl\{(q_C)_{C\in\mathcal C}\;:\;
      &q_C\in\Delta(2^{|C|})\ \text{for each }C;\\
      &q_C\restr S=q_D\restr S\ \text{for every edge }(C,D)\in T,\ S=C\cap D;\\
      &\text{every hosted constraint holds on its host cluster}\bigr\}.
  \end{align*}
  \item A constraint $c$ is \emph{hosted} if some cluster
    $C_h\supseteq\scope(c)$ exists; its row is then imposed on $q_{C_h}$'s
    marginal.
  \item An LMC assertion $(X\indep Y\mid S)$ is \emph{separator-entailed} if $S$
    separates $X$ from $Y$ in $G^\ast$.
  \item For a query atom $a$ and evidence $e$ lying together in a cluster $C_a$,
    the \emph{read-off} is $\readoff_a(q)=\sum_{i:a=1}q_{C_a}[i]$ (no evidence),
    or $\readoff_a(q)=\bigl(\sum_{i:a=1,e}q_{C_a}[i]\bigr)/
    \bigl(\sum_{i:e}q_{C_a}[i]\bigr)$ (with evidence).
  \item $\pi(p)=(p\restr C)_{C\in\mathcal C}$ is the \emph{projection} of a joint
    onto the clusters.
\end{itemize}
\end{definition}

The reconstruction direction needs the following lemma, whose usual statement
contains two defects we must repair: it is normally asserted with a
\emph{uniqueness} claim that is false in the presence of zero separator cells,
and the division in the product formula is left unaddressed.

\begin{lemma}[RIP reconstruction, with zeros]\label{lem:rip}
Let $T$ be a clique tree with the running-intersection property over clusters
$\mathcal C$, and let $(q_C)\in\prod_C\Delta(2^{|C|})$ be separator-consistent:
$q_C\restr S=q_D\restr S$ for every edge $(C,D)$ with $S=C\cap D$. Then there
exists a joint distribution $p\in\Delta(2^{\mathbf V})$ with $p\restr C=q_C$ for
every $C\in\mathcal C$, and $p$ factorizes over the cliques of $G^\ast$. Such $p$
is unique if every separator marginal is strictly positive, but \emph{not} in
general.
\end{lemma}

\begin{proof}
Induction on $|\mathcal C|$. If $|\mathcal C|=1$ take $p=q_C$.

Let $|\mathcal C|\ge2$ and pick a leaf cluster $C$ of $T$ with neighbour $D$ and
separator $S=C\cap D$. Let $T'=T-C$; by the running-intersection property
$C\setminus S$ meets no cluster of $T'$, and $(q_{C'})_{C'\in T'}$ is still
separator-consistent, so the inductive hypothesis gives a joint $p'$ over
$\mathbf V'=\bigcup_{C'\in T'}C'$ with $p'\restr C'=q_{C'}$ for all
$C'\in T'$. Note
$p'\restr S=(p'\restr D)\restr S=q_D\restr S=q_C\restr S=:q_S$, using separator
consistency.

Define a conditional kernel on the fresh atoms $C\setminus S$: for each
assignment $s$ of $S$,
\[
  \kappa(\,\cdot\mid s)=
  \begin{cases}
    q_C(\,\cdot\,,s)/q_S(s), & q_S(s)>0,\\
    \delta_{c_0}, & q_S(s)=0,
  \end{cases}
\]
where $c_0$ is any fixed assignment of $C\setminus S$. This is well defined: when
$q_S(s)>0$ the right-hand side is non-negative and sums over $C\setminus S$ to
$q_S(s)/q_S(s)=1$, so it is a distribution. Now set
$p(x)=p'(x_{\mathbf V'})\cdot\kappa(x_{C\setminus S}\mid x_S)$.

$p$ is a distribution: it is non-negative, and summing over $x_{C\setminus S}$
first gives $p'$, which sums to $1$.

\emph{The arbitrary choice is irrelevant.} On $\{x:q_S(x_S)=0\}$ we have
$p'(x_{\mathbf V'})=0$, because $p'\restr S=q_S$ assigns zero mass to that event
and $p'\ge0$. Hence $p$ vanishes there regardless of $\kappa$, so $p$ does not
depend on $c_0$. (This is also the precise sense in which the familiar quotient
formula $p=\prod_C q_C/\prod_{(C,D)}q_S$ is legitimate: a zero separator cell
forces every contributing cluster cell to zero, since $q_S=q_C\restr S$ is a sum
of non-negative $q_C$ entries, so numerator and denominator vanish together and
the term is assigned $0$.)

\emph{Cluster marginals.} For $C'\in T'$: summing $p$ over $x_{C\setminus S}$
gives $p'$, so $p\restr C'=p'\restr C'=q_{C'}$. For $C$ itself:
$p\restr C(c,s)=\bigl(\sum$ of $p$ over $\mathbf V'\setminus S$ at fixed
$x_S=s\bigr)\cdot\kappa(c\mid s)=p'\restr S(s)\,\kappa(c\mid s)
=q_S(s)\kappa(c\mid s)$, which equals $q_C(c,s)$ when $q_S(s)>0$ and equals
$0=q_C(c,s)$ when $q_S(s)=0$. So $p\restr C=q_C$.

\emph{Factorization.} Unwinding the induction, $p$ is a product of one
non-negative factor per cluster (namely $q_C$ or the kernel $\kappa$ derived from
it, each a function of the atoms of that cluster only), hence factorizes over the
cliques of $G^\ast$.

\emph{Non-uniqueness.} If some separator cell $q_S(s)=0$, then
$\kappa(\cdot\mid s)$ was chosen arbitrarily; any other choice yields a joint
with the same cluster marginals, since all such joints vanish on that event.
Hence $p$ is not unique. When all separator marginals are positive the kernels
are forced and the induction determines $p$ uniquely.
\end{proof}

\begin{theorem}[\textsc{CredalJT} exactness]\label{thm:d5}
Suppose \textup{(H1)} $T$ has the running-intersection property (true by
construction); \textup{(H2)} every LCN sentence is hosted; \textup{(H3)} every
LMC assertion is either hosted or separator-entailed in $G^\ast$; \textup{(H4)}
the query atom and all evidence atoms lie together in some cluster. Then for
every such atom $a$,
\[
  \min_{q\in\mathcal G}\readoff_a(q)=\min_{p\in\DLCN}\readoff_a(\pi(p))=\lo P(a),
  \qquad
  \max_{q\in\mathcal G}\readoff_a(q)=\max_{p\in\DLCN}\readoff_a(\pi(p))=\hi P(a),
\]
and likewise for a conditional query $P(a\mid e)$ subject to
Remark~\ref{rem:denfloor}. Moreover the \emph{projection} half holds without
\textup{(H2)}/\textup{(H3)}, so \textsc{CredalJT} is always a valid \emph{outer}
bound.
\end{theorem}

\begin{proof}
First note the read-off identity: for any joint $p$,
$\readoff_a(\pi(p))=P_p(a)$ (resp.\ $P_p(a\mid e)$), because $\readoff_a$ sums
the cells of $\pi(p)_{C_a}=p\restr C_a$ in which $a$ (and $e$) hold, which by
definition of a marginal equals $P_p(a)$ (resp.\ the ratio). So both sides
optimize the \emph{same} functional.

\emph{Projection $\DLCN\to\mathcal G$ (needs only H1, H4).} Let $p\in\DLCN$ and
set $q_C:=p\restr C$, i.e.\ $q=\pi(p)$. Each $q_C$ is a distribution (a marginal
of one). Each separator equality holds because both $q_C\restr S$ and
$q_D\restr S$ equal $p\restr S$. Each hosted constraint $c$ holds: by
Lemma~\ref{lem:locality} its value depends on $p$ only through
$p\restr\scope(c)$, and since $\scope(c)\subseteq C_h$ we have
$q_{C_h}\restr\scope(c)=p\restr\scope(c)$, so $c$ holds on the host exactly
because it holds for $p$. Hence $q\in\mathcal G$, with
$\readoff_a(q)=P_p(a)$. Therefore every value achievable over $\DLCN$ is
achievable over $\mathcal G$, giving
$\min_{\mathcal G}\readoff_a\le\min_{\DLCN}$ and
$\max_{\mathcal G}\readoff_a\ge\max_{\DLCN}$: the cluster interval
\emph{contains} the true interval. This is the outer-bound claim, and it used
neither H2 nor H3.

\emph{Reconstruction $\mathcal G\to\DLCN$ (needs H2, H3).} Let
$q\in\mathcal G$. By H1 and Lemma~\ref{lem:rip} there is a joint $p$ with
$p\restr C=q_C$ for every cluster, factorizing over the cliques of $G^\ast$. We
check $p\in\DLCN$.

Sub-scope consistency: if $A\subseteq C$ for some cluster $C$, then
$p\restr A=(p\restr C)\restr A=q_C\restr A$, since marginalizing in stages agrees
with marginalizing at once.

\emph{Hosted constraints} (H2 and the hosted part of H3): let $c$ be hosted on
$C_h$. By the previous paragraph $p\restr\scope(c)=q_{C_h}\restr\scope(c)$, and
$q$ satisfies $c$ on $C_h$ by definition of $\mathcal G$; by
Lemma~\ref{lem:locality} $c$ therefore holds for $p$.

\emph{Separator-entailed LMC} (the other part of H3): let
$(X\indep Y\mid S)$ be separator-entailed, so $S$ separates $X$ from $Y$ in
$G^\ast$. Since $p$ factorizes over the cliques of $G^\ast$, it satisfies the
\emph{global Markov property} of $G^\ast$: factorization implies global Markov
for any undirected graph, and crucially this direction requires \emph{no}
positivity assumption (Lauritzen, \emph{Graphical Models}, Prop.~3.8; only the
converse, global Markov $\Rightarrow$ factorization, needs $p>0$ via
Hammersley--Clifford). Global Markov states that if $S$ separates $A$ from $B$ in
$G^\ast$ then $A\indep B\mid S$ under $p$; applying it to the connected
components separated by $S$ and then restricting to the subsets $X$ and $Y$ of
those components (conditional independence is preserved under taking subsets, by
the decomposition property of conditional independence) gives
$X\indep Y\mid S$ under $p$, i.e.\ \eqref{eq:lmc} holds. Note this assertion was
never imposed on $\mathcal G$: it holds automatically for the reconstructed $p$.

Hence $p\in\DLCN$, and $\readoff_a(q)=\readoff_a(\pi(p))=P_p(a)$ because
$\pi(p)=q$ on the host cluster. So every value achievable over $\mathcal G$ is
achievable over $\DLCN$, giving the reverse inequalities. Combined with the
projection half, the optima coincide. For a conditional query, by H4 the read-off
is a functional of the single cluster $C_a$ and both maps preserve $q_{C_a}$, so
the same argument applies verbatim.
\end{proof}

\begin{remark}[Reading the proof]\label{rem:d5reading}
(a) The argument never mentions chains: exactness is a property of the
\emph{junction tree}, via RIP, hosting, separator-entailment and locality. It is
in particular \emph{not} a statement about singly-connected topologies --- the
worked example of Section~\ref{sec:example} is loopy and exact.
(b) The projection half is unconditional, so \textsc{CredalJT} is \emph{never
unsound}: its interval always contains the truth, and H2/H3 decide only
tightness. (c) H3 is \emph{sufficient, not necessary}: on
\texttt{examples/asia.lcn} and \texttt{examples/polytree.lcn} some assertions are
neither hosted nor separator-entailed, yet \textsc{CredalJT} is empirically exact
--- the unentailed assertions simply do not bind. Where H3 fails \emph{and}
binds, the result is a valid but loose outer bound: on
\texttt{examples/smokers.lcn} the six large-scope LMCs of the undirected
component $F_1F_2F_3$ are unhostable and unentailed (within the width-$6$ cluster
$\{F_1,F_2,F_3,S_1,S_2,S_3\}$ no separator divides $S_3$ from $F_1$), and
\textsc{CredalJT} returns $P(F_1)=[0,1]$.
\end{remark}

\begin{remark}[The denominator floor is a soundness hypothesis]\label{rem:denfloor}
The conditional read-off is linearized by the auxiliary equality
\eqref{eq:nlp-obj-ev} together with the denominator floor \eqref{eq:nlp-floor}.
The floor is an extra constraint on $\mathcal G$ that $\DLCN$ does not have, so
the projection half of Theorem~\ref{thm:d5} --- and with it the outer-bound
guarantee --- holds only for joints with $P(e)\ge\varepsilon$. If
$\min_{p\in\DLCN}P(e)<\varepsilon$, the returned conditional interval can be
\emph{unsound}. This is a soundness caveat, not merely a tightness one, and it is
the price of admitting conditional queries into a globally solved formulation.
\end{remark}

\begin{corollary}[Exactness on \texttt{ktree-fr} at treewidth cost]\label{cor:ktree}
On every \texttt{ktree-fr} instance, H1--H4 hold, so \textsc{CredalJT} returns
the exact $\DLCN$ marginal bounds at cost $O(n\,2^{k+1})$.
\end{corollary}

\begin{proof}
A $k$-tree is chordal, and in the \texttt{ktree-fr} class it is oriented as a DAG
in construction order, every atom conditioning on the \emph{full conjunction} of
its $k$ clique-parents. Those parents are pairwise adjacent in the $k$-tree, so
moralization adds no edge and triangulation is a no-op: $G^\ast$ is the $k$-tree
itself and its maximal cliques --- the clusters --- are exactly the
$(k{+}1)$-cliques, giving H1 by construction. Every sentence is either a root
marginal or scoped to a child with its $\le k$ parents, which co-occur in a
cluster: H2. Each LMC assertion produced by the Local Markov Condition has as its
conditioning set a full clique-separator of the $k$-tree, which separates the two
sides in $G^\ast$, so each is hosted or separator-entailed: H3. Any single query
atom lies in a $(k{+}1)$-clique: H4. Theorem~\ref{thm:d5} applies, and the cost
is $\sum_C2^{|C|}=O(n)\cdot2^{k+1}$, over $O(n)$ clusters each of size $k{+}1$.
\end{proof}

\begin{remark}[What the \texttt{ktree-fr} corollary does and does not claim]\label{rem:ktreehonest}
Corollary~\ref{cor:ktree} is a statement about \emph{cost}: exactness at
$O(n2^{k+1})$ rather than $2^n$. It should \emph{not} be read as ``the one loopy
topology where a treewidth-bounded engine beats the strong-extension engines on
accuracy''. On the loopy realizable $k$-tree of Section~\ref{sec:example} Credal
VE is also exact, atom for atom. On a larger $k{=}3$ instance
(\texttt{ktree\_fr\_k3\_n5\_1.lcn}) Credal VE fails not by being loose but by not
finishing: the credal network there is tiny ($512$ vertex combinations across
five families), yet elimination did not terminate within $200$\,s, stalling
inside potential pruning, while \textsc{CredalJT} solved the instance in about a
second. The separation between \textsc{CredalJT} and Credal VE on this family is
therefore one of \emph{tractability}, not of accuracy. Loopiness alone is not a
reason to expect looseness.
\end{remark}

\begin{example}[Verified end to end]\label{ex:d5verified}
Section~\ref{sec:example} is the worked case at $k=2$. Independently:
on \texttt{ktree\_fr\_k3\_n5\_1.lcn} the two maximal cliques are
$\{x_0,x_1,x_3,x_4\}$ and $\{x_1,x_2,x_3,x_4\}$ (treewidth $3$), the sole LMC
assertion $(x_0\indep x_2\mid x_1,x_3,x_4)$ has conditioning set equal to the
separator between them, and \textsc{CredalJT} reproduces the certified
certified global $2^n$ bounds on all five atoms; on
\texttt{examples/chain.lcn} the clusters are the $7$ edge-pairs (treewidth $1$:
$28$ program variables against a $2^8=256$ joint); and on
\texttt{examples/polytree.lcn} the collider cliques $\{x_0,x_1,x_2\}$ and
$\{x_3,x_4,x_5\}$ give treewidth $2$.
\end{example}

\subsection{Position among the engines}\label{sec:theory-position}

Table~\ref{tab:master} places \textsc{CredalJT} against the other shipped
engines, and Table~\ref{tab:topo} reports the observed bound character by
topology.

\begin{table}[ht]\centering\footnotesize
\caption{Bound character per engine. ``Targets'' is the feasible set the engine
optimizes; ``character'' is per Definition~\ref{def:char} \emph{against that
set}. The fourth column gives the condition for the returned interval to equal
the $\DLCN$ truth.}\label{tab:master}
\setlength{\tabcolsep}{4pt}
\begin{tabular}{@{}lllll@{}}
\toprule
Engine & Targets & Character (vs.\ target) & Exact w.r.t.\ $\DLCN$ when & Cost \\
\midrule
Direct $2^n$, global & $\DLCN$ & exact if certified &
  endpoint certified & $2^n$, global \\
Direct $2^n$, local  & $\DLCN$ & unreliable &
  --- (may over-tighten) & $2^n$, ipopt \\
\textbf{\textsc{CredalJT}} (D5)  & $\DLCN$ & outer always; exact if H1--H4 &
  H1--H4 (Thm.~\ref{thm:d5}) & $O(n2^{w+1})$ \\
\midrule
Credal~VE                        & $\DSE$ & exact (unconditional queries) &
  no Gap~A and no Gap~B & per-target elim.\ \\
Credal~CTE ($\epsilon{=}0$)      & $\DSE$ & \textbf{inner} &
  clean s.c.; no cross-family sent. & two-pass tree \\
ApproxLP                         & $\DSE$ & \emph{inner} &
  no gaps $+$ no stall & coord.\ descent \\
Interval~BP (interval)  & $\DSE$ & outer; exact on tree &
  tree $+$ no gaps & loopy messages \\
Credal~IJGP                      & $\DSE$ & outer; exact if $i\ge w$ &
  $i\ge w$ $+$ no gaps & $2^i$/cluster \\
\midrule
ARIEL                            & $\DLCN$ & outer always &
  see the ARIEL note & factor messages \\
\bottomrule
\end{tabular}
\end{table}

\begin{table}[ht]\centering\footnotesize
\caption{Bound character by topology, unconditional marginals. ``ex'' $=$ exact
w.r.t.\ $\DLCN$, ``out'' $=$ valid outer bound, ``in'' $=$ inner (unsound as a
guarantee), ``n/f'' $=$ did not finish.}\label{tab:topo}
\setlength{\tabcolsep}{4.5pt}
\begin{tabular}{@{}lcccccc@{}}
\toprule
Topology & Exact-glob & \textsc{CredalJT} & CVE & CTE & ApproxLP & IBP-int \\
\midrule
Markov chain, tree, polytree (realizable) & ex & ex & ex & ex & ex & ex \\
Chain graph with blocks (\texttt{alarm})  & ex$^\ast$ & ex & ex & ex & ex & ex \\
Cross-family sentence (\texttt{d4\_biting}) & ex & ex & ex & \textbf{in} & ex & ex \\
Non-realizable Type-1                     & ex & ex & out & out & out & out \\
Non-realizable Type-2                     & ex & ex & out & out & out & out \\
Loopy $k$-tree, realizable, $k{=}2$ (\S\ref{sec:example}) & ex & ex & \textbf{ex} & --- & --- & --- \\
Loopy $k$-tree, realizable, $k{=}3$       & ex & ex & \textbf{n/f} & --- & --- & --- \\
Undirected block (\texttt{smokers})       & ex & \textbf{out} & out & out & in & out \\
\bottomrule
\end{tabular}
\end{table}

\noindent$^\ast$ On \texttt{alarm} the global backend leaves several endpoints
uncertified and the \emph{local} variant over-tightens; the certified
values come from the dedicated verifiers.

In summary: on bounded-treewidth instances \textsc{CredalJT} is exact under
H1--H4 and a valid outer bound otherwise, at $O(n2^{w+1})$ rather than $2^n$; it
remains applicable on instances where the direct formulation is solvable but its
global solver certifies no endpoint. For conditional queries
Remark~\ref{rem:denfloor} applies.

\section{A self-contained subsection for a paper}\label{sec:paper}

% ======================================================================
% BEGIN LIFTABLE BLOCK -- self-contained; no dependence on the rest of
% this document. Renumber the theorem and adjust the appendix pointer.
% ======================================================================

\subsection*{Exact inference at treewidth cost}

Let an LCN over binary atoms $\mathbf V$ consist of sentences bounding
$P(\varphi)$ or $P(\varphi\mid\psi)$, read in cleared form
$\ell P(\psi)\le P(\varphi\wedge\psi)\le uP(\psi)$, together with the
independence assertions $(X\indep Y\mid S)$ derived by the Local Markov
Condition, imposed as $P(\mathbf x,\mathbf y,\mathbf s)P(\mathbf s)
=P(\mathbf x,\mathbf s)P(\mathbf y,\mathbf s)$. These determine a credal set
$\DLCN\subseteq\Delta(2^{\mathbf V})$, and inference seeks
$\lo P(a)=\min_{\DLCN}P(a)$ and $\hi P(a)=\max_{\DLCN}P(a)$. Optimizing over
$\DLCN$ directly costs $2^n$ variables; compiling to a credal network and
propagating is cheaper but optimizes a strictly larger set, the strong extension.

\textsc{CredalJT} avoids both. Moralize the chain graph induced by the program's
families, triangulate, and build a clique tree $T$ over clusters $\mathcal C$
with the running-intersection property; let $G^\ast$ be the associated chordal
graph. Give each cluster a variable vector $q_C\in\Delta(2^{|C|})$ over the joint
of its atoms and solve
\[
\begin{aligned}
\min/\max\quad & \readoff_a(q)\\
\text{s.t.}\quad
&\textstyle\sum_i q_{C,i}=1,\ q_{C,i}\ge0 && \forall C\in\mathcal C,\\
&q_C\restr S=q_D\restr S && \forall (C,D)\in T,\ S=C\cap D,\\
&\ell\langle A_{\psi},q_{C_h}\rangle\le\langle A_{\varphi\wedge\psi},q_{C_h}\rangle
   \le u\langle A_{\psi},q_{C_h}\rangle && \forall\text{ sentence hosted on }C_h,\\
&\langle A_a,q\rangle\langle A_b,q\rangle=\langle A_c,q\rangle\langle A_d,q\rangle
   && \forall\text{ assertion hosted,}
\end{aligned}
\]
where $A_\varphi$ is the $0/1$ indicator of $\varphi$ over the host cluster's
states, a constraint is \emph{hosted} on any cluster containing its scope, and
$\readoff_a$ reads $P(a)$ (or $P(a\mid e)$, via one auxiliary variable) off a
cluster containing $a$ and the evidence. An assertion whose conditioning set
separates its two sides in $G^\ast$ is \emph{separator-entailed} and is omitted:
the tree's topology already enforces it. The program is nonconvex through the
bilinear assertion rows and is solved globally.

\begin{theorem}\label{thm:paper}
Assume \textup{(H1)} $T$ has the running-intersection property; \textup{(H2)}
every sentence is hosted; \textup{(H3)} every assertion is hosted or
separator-entailed in $G^\ast$; \textup{(H4)} the query and evidence atoms share
a cluster. Then the program's optima equal $\lo P(a)$ and $\hi P(a)$ exactly.
Moreover the direction giving containment of the true interval requires only
\textup{H1} and \textup{H4}, so the returned interval is always a valid outer
bound.
\end{theorem}

The proof has two halves. Projection: for $p\in\DLCN$, the cluster marginals
$p\restr C$ are separator-consistent and satisfy every hosted row, because a
constraint's value depends on $p$ only through its scope, which lies inside the
host; hence every value attained on $\DLCN$ is attained on the relaxation, which
is the outer-bound claim. Reconstruction: a separator-consistent family of
cluster marginals reassembles, by running intersection, into a joint $p$ with
those marginals which factorizes over the cliques of $G^\ast$ (the reassembly
needs care where separator cells vanish: the conditional kernels are then
arbitrary but immaterial, since $p$ vanishes there regardless). Hosted rows
transfer back to $p$ by locality, and separator-entailed assertions hold
automatically, since factorization implies the global Markov property of
$G^\ast$ with no positivity assumption. Full proofs are in the appendix.

The cost is $\sum_C2^{|C|}=O(n2^{w+1})$ variables in the treewidth $w$, against
$2^n$ for the direct formulation, plus $O(n2^w)$ separator equalities. The
constraint system is query-independent, so all marginals are obtained from one
program by swapping the objective. As a consequence, on a $k$-tree rendered so
that each atom conditions on the full conjunction of its $k$ clique-parents,
moralization and triangulation are no-ops, every hypothesis holds, and exact
bounds are obtained at $O(n2^{k+1})$. Empirically, on a loopy four-atom $k$-tree
the formulation reproduces a certified global $2^n$ reference on every atom,
using $22$ variables and $22$ constraints.

% ======================================================================
% END LIFTABLE BLOCK
% ======================================================================

\section{Experimental setup}\label{sec:repro}

All reported bounds were computed on the stated instances with the cluster program
of Section~\ref{sec:nlp}, solved by a spatial branch-and-bound global solver at
gap tolerance $0$, and independently cross-checked against the direct $2^n$
formulation over $\DLCN$ under the same global solver. On the instance of
Section~\ref{sec:example} the reference reported status \emph{confirmed} at
optimality gap $0$ for all eight endpoint optimizations, and the two formulations
agreed on every atom to the precision displayed.

The junction trees, separator equalities and constraint counts of
Section~\ref{sec:example} and Table~\ref{tab:contrast} were read off the
constructed cluster tree and constraint system directly rather than reconstructed
by hand. All instances are shipped with the library: the running example is a
$k$-tree with $k=2$ on four atoms, generated in the family-realizable regime, and
the two instances of Table~\ref{tab:contrast} are the standard \texttt{asia} and
\texttt{alarm} programs.

\section{Provenance}\label{sec:provenance}

Every interval and every program-size figure in this document was computed as
described in Section~\ref{sec:repro}. A value attributed to the global $2^n$
reference is reported as exact only where that optimization was certified at gap
$0$.

Two classes of statement are cited rather than re-derived here. The
\texttt{smokers} and \texttt{asia}/\texttt{polytree} observations of
Remark~\ref{rem:d5reading}, the \texttt{alarm} certification caveat under
Table~\ref{tab:topo}, and the comparative rows of Tables~\ref{tab:master}
and~\ref{tab:topo} are taken from the consolidated reference
\texttt{docs/properties.tex}, which reports them with its own provenance. The
non-termination of Credal VE on the $k{=}3$ instance in
Remark~\ref{rem:ktreehonest} is likewise cited from that document.

Related notes in \texttt{docs/}: \texttt{properties.tex} consolidates the
exactness results for all engines and is the source of the theory reproduced in
Section~\ref{sec:theory}; \texttt{cjt\_exactness.tex} is a forensic companion
investigating a reported \textsc{CredalJT}-versus-oracle discrepancy on
\texttt{ktree\_fr\_k3\_n5\_1.lcn} that does not reproduce;
\texttt{tighter\_approximation.tex} defines the D1--D5 scheme taxonomy in which
\textsc{CredalJT} is D5; and \texttt{lrs.tex} covers the vertex enumeration that
\textsc{CredalJT}, uniquely among the credal-network engines, does not need.

\end{document}
